%=====================================================================
\chapter{Statistical Test Selection and Formula Sheet}\label{app:stats}
%=====================================================================

A working reference for Part~IV. Nothing here replaces understanding what the
test assumes (\chref{inference}) --- it is here so you do not have to remember
which function to call.

\section{Choosing a Test}

Read down the first two columns to your situation.

\begin{center}
\small
\begin{longtable}{@{}L{3.2cm}L{2.6cm}L{4.0cm}L{4.2cm}@{}}
\toprule
\textbf{Question} & \textbf{Outcome} & \textbf{Parametric} & \textbf{Robust alternative} \\
\midrule
\endfirsthead
\toprule
\textbf{Question} & \textbf{Outcome} & \textbf{Parametric} & \textbf{Robust alternative} \\
\midrule
\endhead
\bottomrule
\endfoot
One group vs a value
  & Continuous
  & One-sample $t$
  & Wilcoxon signed-rank; bootstrap CI \\
Two independent groups
  & Continuous
  & Welch's $t$ (\emph{not} Student's)
  & Mann--Whitney $U$; bootstrap \\
Two paired measurements
  & Continuous
  & Paired $t$
  & Wilcoxon signed-rank \\
$k$ independent groups
  & Continuous
  & One-way ANOVA
  & Kruskal--Wallis \\
$k$ paired measurements
  & Continuous
  & Repeated-measures ANOVA
  & Friedman \\
Two factors
  & Continuous
  & Two-way ANOVA
  & Aligned rank transform ANOVA \\
Nested / repeated data
  & Continuous
  & Linear mixed model (\chref{regression})
  & Same, with bootstrap CIs \\
Association
  & Both continuous
  & Pearson $r$
  & Spearman $\rho$; Kendall $\tau$ \\
Association
  & Both categorical
  & Chi-square test of independence
  & Fisher's exact (small expected counts) \\
Agreement, 2 raters
  & Categorical
  & ---
  & Cohen's $\kappa$ \\
Agreement, $n$ raters
  & Any
  & ---
  & Krippendorff's $\alpha$ \\
Predict a continuous outcome
  & Continuous
  & Linear regression
  & Robust regression; quantile regression \\
Predict a binary outcome
  & Binary
  & Logistic regression
  & --- \\
Predict a count
  & Count
  & Poisson --- only if variance $\approx$ mean
  & Negative binomial (usually correct for software data) \\
Time to event
  & Duration
  & Cox proportional hazards
  & Kaplan--Meier with log-rank \\
Classifiers over datasets
  & Performance
  & ---
  & Friedman + Nemenyi (\chref{mleval}) \\
``No meaningful difference''
  & Continuous
  & TOST equivalence test
  & Bayes factor \\
\end{longtable}
\end{center}

\begin{alertbox}[Before any of them]
Plot the data. A boxplot with the points overlaid, and a Q--Q plot of the
residuals, will tell you more in ten seconds than the test will --- including
whether the test is appropriate at all. Testing for normality is not the way to
decide (\chref{inference}).
\end{alertbox}

\section{Formulas Worth Having}

\subsection*{Descriptive}

\[
\bar{x} = \frac{1}{n}\sum_{i=1}^{n} x_i
\qquad
s = \sqrt{\frac{1}{n-1}\sum_{i=1}^{n}(x_i - \bar{x})^2}
\qquad
\mathrm{SE} = \frac{s}{\sqrt{n}}
\]

\subsection*{Confidence interval for a mean}

\[
\bar{x} \pm t_{\alpha/2,\;n-1}\cdot \frac{s}{\sqrt{n}}
\]

\noindent A 95\% interval means: if the study were repeated many times, 95\% of
such intervals would contain the true value. It does \emph{not} mean there is a
95\% probability the true value lies in this one (\chref{inference}).

\subsection*{Effect sizes}

\[
d = \frac{\bar{x}_1 - \bar{x}_2}{s_{\text{pooled}}},
\qquad
s_{\text{pooled}} = \sqrt{\frac{(n_1-1)s_1^2 + (n_2-1)s_2^2}{n_1+n_2-2}}
\]

\[
g = d\left(1 - \frac{3}{4(n_1+n_2)-9}\right)
\qquad\text{(Hedges' correction; use below about $n=20$ per group)}
\]

\[
\delta = \frac{\#\{x_i > y_j\} - \#\{x_i < y_j\}}{n_1 n_2}
\qquad\text{(Cliff's $\delta$; distribution-free, $-1$ to $1$)}
\]

\[
\text{OR} = \frac{p_1/(1-p_1)}{p_2/(1-p_2)}
\qquad
\text{RR} = \frac{p_1}{p_2}
\]

\noindent Report the risk ratio alongside the odds ratio: they diverge sharply
when the outcome is common, and the odds ratio always looks larger
(\chref{regression}).

\subsection*{Agreement}

\[
\kappa = \frac{p_o - p_e}{1 - p_e}
\]

\noindent $p_o$ observed agreement, $p_e$ agreement expected by chance from the
observed marginals (\chref{slr}, \chref{qualitative}).

\subsection*{Multiple comparisons}

\begin{center}
\small
\begin{tabular}{@{}L{3.6cm}L{10.7cm}@{}}
\toprule
Bonferroni  & Reject $H_i$ if $p_i < \alpha/m$ \\
Holm        & Sort ascending; reject while $p_{(i)} < \alpha/(m-i+1)$; stop at
              the first failure. Uniformly better than Bonferroni \\
Benjamini--Hochberg & Sort ascending; find the largest $k$ with
              $p_{(k)} \le \frac{k}{m}\alpha$; reject all up to $k$ \\
\bottomrule
\end{tabular}
\end{center}

\subsection*{Interpreting effect sizes}

\begin{center}
\small
\begin{tabular}{@{}L{3.0cm}L{2.6cm}L{2.6cm}L{2.6cm}L{2.4cm}@{}}
\toprule
 & \textbf{Negligible} & \textbf{Small} & \textbf{Medium} & \textbf{Large} \\
\midrule
Cohen's $d$      & $<0.2$ & 0.2 & 0.5 & 0.8 \\
Cliff's $\delta$ & $<0.147$ & 0.147 & 0.33 & 0.474 \\
Pearson $r$      & $<0.1$ & 0.1 & 0.3 & 0.5 \\
\bottomrule
\end{tabular}
\end{center}

\begin{alertbox}[These bands are conventions, not findings]
Cohen proposed them for fields with no prior information and warned repeatedly
against routine use. A $d$ of 0.2 may be practically enormous where the outcome
is costly and the intervention is free, and a $d$ of 0.9 may be worthless if
the intervention cannot be deployed.

\medskip
Interpret against your smallest effect size of interest (\chref{sampling}), not
against this table. The table is here because reviewers ask for it.
\end{alertbox}

\section{Sample Size Quick Reference}

Approximate $n$ \emph{per group}, two-sided $\alpha = 0.05$, power 0.80,
two-group comparison of means:

\begin{center}
\small
\begin{tabular}{@{}L{3.0cm}ccccc@{}}
\toprule
Effect $d$        & 0.20 & 0.35 & 0.50 & 0.80 & 1.00 \\
\midrule
$n$ per group     & 394  & 130  & 64   & 26   & 17 \\
Total $N$         & 788  & 260  & 128  & 52   & 34 \\
\bottomrule
\end{tabular}
\end{center}

\noindent Within-subjects designs need substantially fewer participants, by a
factor depending on the correlation between conditions --- which is the main
practical reason to prefer them (\chref{experiments}). Always compute your own
with \texttt{pwr} rather than reading off this table; it is here for a sense of
scale.

\section{Code}

Runnable versions of everything above are in the \texttt{code/} directory
alongside this handout, keyed by chapter.

\begin{lstlisting}[language=Rlang, caption={\texttt{code/appB\_tests.R} --- the
common cases, with the reporting each needs.}]
# --- Two independent groups -------------------------------------------
t.test(y ~ g, data = d)                      # Welch by default in R
effsize::cohen.d(y ~ g, data = d)            # with CI
effsize::cliff.delta(y ~ g, data = d)        # distribution-free

# --- Paired ------------------------------------------------------------
t.test(pre, post, paired = TRUE)
wilcox.test(pre, post, paired = TRUE, conf.int = TRUE)

# --- Nested / repeated -------------------------------------------------
m <- lme4::lmer(y ~ cond + (1 + cond | subject) + (1 | item), data = d)
confint(m, method = "profile")

# --- Counts ------------------------------------------------------------
m <- MASS::glm.nb(defects ~ size + churn, data = d)   # not glm(poisson)
performance::check_overdispersion(m)

# --- Power -------------------------------------------------------------
pwr::pwr.t.test(d = 0.35, sig.level = 0.05, power = 0.80,
                type = "two.sample")

# --- Equivalence -------------------------------------------------------
TOSTER::tsum_TOST(m1 = 6.1, m2 = 5.9, sd1 = 2.4, sd2 = 2.6,
                  n1 = 60, n2 = 60, eqb = 0.40)

# --- Multiplicity ------------------------------------------------------
p.adjust(pvals, method = "holm")   # confirmatory
p.adjust(pvals, method = "BH")     # exploratory
\end{lstlisting}

\begin{lstlisting}[language=PythonARM, caption={\texttt{code/appB\_tests.py} --- the
same cases in Python, for those who prefer it.}]
import numpy as np, pingouin as pg, statsmodels.formula.api as smf

pg.ttest(a, b, correction=True)          # Welch
pg.compute_effsize(a, b, eftype="cohen")
pg.mwu(a, b)                             # Mann-Whitney

# Mixed model
smf.mixedlm("y ~ cond", data=d, groups=d["subject"]).fit().summary()

# Negative binomial for counts
smf.negativebinomial("defects ~ size + churn", data=d).fit().summary()

# Power
from statsmodels.stats.power import TTestIndPower
TTestIndPower().solve_power(effect_size=0.35, alpha=0.05, power=0.80)

# Multiplicity
from statsmodels.stats.multitest import multipletests
multipletests(pvals, method="holm")
\end{lstlisting}

\begin{conceptbox}[Reporting checklist for any quantitative result]
\begin{tight}
  \item The estimate, in natural units.
  \item A confidence interval on it.
  \item An effect size, with its interval.
  \item $n$, at every level of nesting.
  \item The test used, and why it rather than the alternative.
  \item Assumption checks performed, and what they showed.
  \item Any multiplicity correction, and which comparisons it covered.
  \item The p-value, last, as one piece of information among several.
\end{tight}
\end{conceptbox}
